%% kapa: exam=2 round=33 year=2022 period=1 subject=appraisal_practice subject_ko=감정평가실무 seq=2 number=2 points=30
%% source: Q-Net 감정평가사 기출문제 (cst00309) artlSeq=5210838 file=raw/1008/5210838/2022년 제33회 감정평가사 2차시험 1교시 문제지(감정평가실무).pdf
%% conversion: kordoc 4.24.51 -> tools/gen2.py (auto); review=proofread-gpt-astra; flags=none
\begin{longquestion}{【문제 2】}
평가대상 토지는 주거지역의 소로에 접한 나지이다. 감정평가사 甲은 거래사례비교법으로 대상 토지를 평가하기 위하여 인근의 주거지역에서 최근 3년 이내에 거래된 총 4건의 거래관련 사례 등을 아래와 같이 수집하였다. 이들 사례 등에 대하여 거래 당시의 조건 등에 따른 차이를 정상화하기 위한 보정(Adjustments)작업을 한 후, 보정된 토지만의 가격과 단가를 구하고자 한다. (다만, 모든 사례 등은 거래시점과 기준시점의 차이에 따른 가격의 변동은 포착되지 아니하여, 고려하지 아니하기로 함) (30점) \par
{\small\centering\renewcommand{\arraystretch}{1.25}
\resizebox{\ifdim\width>\linewidth\linewidth\else\width\fi}{!}{
\begin{tabular}{|c|c|c|c|c|c|}
\hline
\makecell{구분} & \makecell{거래\\금액 등} & \makecell{토지} & \makecell{건물} & \makecell{거래상황 및 조건} & \makecell{비고} \\
\hline
\makecell{거래사례\\(1)} & \makecell{14억원} & \makecell{500m\textsuperscript{2}} & \makecell{연면적\\1,000m\textsuperscript{2}} & \makecell{공사 중단된 건물을\\포함하여 인수} &  \\
\hline
\makecell{거래사례\\(2)} & \makecell{10억원} & \makecell{500m\textsuperscript{2}} & \makecell{없음} & \makecell{저당대부액 신규인수} &  \\
\hline
\makecell{낙찰사례\\(3)} & \makecell{5억원} & \makecell{400m\textsuperscript{2}} & \makecell{없음} & \makecell{경매에서\\수회 유찰 후 낙찰} & \makecell{점유자 및 적치물\\소재} \\
\hline
\makecell{평가선례\\(4)} & \makecell{10억원} & \makecell{400m\textsuperscript{2}} & \makecell{없음} & \makecell{낙찰사례(3)의\\평가선례} & \makecell{경매목적\\평가} \\
\hline
\end{tabular}}\par}

\medskip\noindent 물음 1) 상기의 거래사례(1)에 대하여, < 자료 1 > 과 < 자료 2 >를 활용, 공사 중단된 상태의 건물의 공정률을 추정하고, 토지만의 보정된 거래가격과 단가를 구하시오. (10점) \par
\medskip\noindent 물음 2) 상기의 거래사례(2)에 대하여, < 자료 3 >과 < 자료 4 >를 활용, 저당대부액의 신규 인수조건을 고려한 저당지불액 및 저당지불액의 현가합 (백원단위에서 반올림)을 구한 후, 토지의 보정된 거래가격과 단가를 구하시오. (10점) \par
\medskip\noindent 물음 3) 상기의 낙찰사례(3)에 대하여, < 자료 5 >를 활용, 토지의 보정된 가격과 단가를 구하되, 낙찰자 입장의 보정된 가격과 평가선례(4)를 활용한 통계적 측면의 보정된 가격을 각각 구한 후, 그 결과를 평균하여 가격과 단가를 구하시오. 또한, 낙찰사례(3)과 평가선례(4)에 대한 가격의 성격 및 보정방식별 특징을 쓰시오. (10점) \par
< 자료 1 > 거래사례(1)에 대한 개요 \par
1. 토지: 면적은 500m\textsuperscript{2}이고, 토지만의 정상화 보정을 위하여 건물은 공정률을 고려한 원가법으로 평가하고, 전체 거래가격에서 공제하는 방식을 채택하여 산정하기로 하였다. \par
2. 건물: 철근콘크리트조 슬래브지붕 3층규모의 원룸형 도시형생활주택으로, 연면적은 1,000m\textsuperscript{2}이고, 최근 공사가 중단된 상태이며, 현상태의 건물과 토지가 총 14억원에 거래되었다. \par
3. 공사진행 사항 \par
- 공정률은 최근 발행된 < 자료 2 >의 건물신축단가표를 활용하여 간접법으로 산정하며, 신축단가(백원단위에서 반올림)로도 활용할 예정이다. \par
- 각 공정별로 철근콘크리트공사, 조적공사, 방수공사 등은 완료되었으나, 대부분의 내부공사가 남은 상태이고, 가설물은 존치된 상태이어서 이를 고려할 때, 가설공사가 90\% 진행된 것으로 추정하였다. \par
- 기초 및 토공사는 건물 주변의 되메우기 및 정리가 남은 상태로 80\% 진행된 것으로 추정하였다. \par
- 이외의 공사는 착수되지 아니하였고, 설비부문에 대한 별도의 보정은 고려치 아니하기로 하였다. \par
- 설계비는 모두 지급되었고, 감리비는 50\%만 지급된 상태이다. \par
- 제경비의 처리에 대하여서는 직접적인 공사비(1～12번 공사)의 집행정도에 따라 그 비율에 의거 추정할 예정이다. \par
< 자료 2 > 건물신축단가표: 공사항목에 따른 법정비율은 고려치 아니한다. \par
1. 개요 \par
{\small\centering\renewcommand{\arraystretch}{1.25}
\resizebox{\ifdim\width>\linewidth\linewidth\else\width\fi}{!}{
\begin{tabular}{|c|c|c|c|}
\hline
\makecell{용도} & \makecell{구조} & \makecell{표준단가(m\textsuperscript{2})} & \makecell{내용연수(년)} \\
\hline
\makecell{도시형생활주택\\(원룸형)} & \makecell{철근콘크리트조\\슬래브지붕} & \makecell{1,600,000} & \makecell{50} \\
\hline
\end{tabular}}\par}

2. 단위면적(m\textsuperscript{2})당 공사비 적산표 \par
{\small\centering\renewcommand{\arraystretch}{1.25}
\resizebox{\ifdim\width>\linewidth\linewidth\else\width\fi}{!}{
\begin{tabular}{|c|c|c|c|}
\hline
\makecell{구분} & \makecell{주요공사내역} & \makecell{공사비} & \makecell{구성비} \\
\hline
\makecell{01. 가설공사} & \makecell{공통가설, 일반가설} & \makecell{80,000} & \makecell{5.00} \\
\hline
\makecell{02. 기초 및 토공사} & \makecell{터파기, 되메우기, 잔토처리, 잡석다짐} & \makecell{20,000} & \makecell{1.25} \\
\hline
\makecell{03. 철근콘크리트공사} & \makecell{레미콘, 철근가공조립, 합판거푸집, 유로폼} & \makecell{280,000} & \makecell{17.50} \\
\hline
\makecell{04. 조적공사} & \makecell{시멘트벽돌쌓기(0.5B), 치장벽돌쌓기\\(0.5B)} & \makecell{20,000} & \makecell{1.25} \\
\hline
\makecell{05. 방수공사} & \makecell{도막방수, 우레탄방수} & \makecell{40,000} & \makecell{2.50} \\
\hline
\makecell{06. 미장공사} & \makecell{시멘트모르타르미장} & \makecell{60,000} & \makecell{3.75} \\
\hline
\makecell{07. 타일공사} & \makecell{자기질타일, 포슬란타일, 석재타일} & \makecell{60,000} & \makecell{3.75} \\
\hline
\makecell{08. 창호공사} & \makecell{칼라알루미늄단열바, 강화도어 등} & \makecell{180,000} & \makecell{11.25} \\
\hline
\makecell{09. 유리공사} & \makecell{로이복층유리24mm,\\강화유리10mm} & \makecell{20,000} & \makecell{1.25} \\
\hline
\makecell{10. 도장공사} & \makecell{페인트 및 수성페인트(내부), 무늬코트} & \makecell{220,000} & \makecell{13.75} \\
\hline
\makecell{11. 수장공사} & \makecell{압출보온판, 걸레받이, 몰딩, 원목마루} & \makecell{200,000} & \makecell{12.50} \\
\hline
\makecell{12. 기타공사} & \makecell{우편함, 카스토퍼, 맨홀} & \makecell{60,000} & \makecell{3.75} \\
\hline
\makecell{소계} &  & \makecell{1,240,000} & \makecell{(77.50)} \\
\hline
\makecell{제경비} & \makecell{간접노무비, 산재보험료, 안전관리비,\\경비, 일반관리비, 이윤, 건강보험료,\\환경보존비} & \makecell{260,000} & \makecell{16.25} \\
\hline
\makecell{건축 공사비 합계} &  & \makecell{1,500,000} & \makecell{(93.75)} \\
\hline
\makecell{설계비} &  & \makecell{24,000} & \makecell{1.50} \\
\hline
\makecell{감리비} &  & \makecell{16,000} & \makecell{1.00} \\
\hline
\makecell{전기기본설비비} &  & \makecell{60,000} & \makecell{3.75} \\
\hline
\makecell{합계} &  & \makecell{1,600,000} & \makecell{100\%} \\
\hline
\end{tabular}}\par}

< 자료 3 > 거래사례(2)에 대한 개요 \par
1. 토지: 면적 500m\textsuperscript{2}이고, 유리한 조건의 신규 저당대부가 설정되어 이의 인수를 전제로 거래가 이루어졌다. \par
2. 건물: 없음 \par
3. 저당대부 관련사항 \par
- 본 거래사례(2)는 매수인이 시장이자율보다 낮은 이율로 승계 가능한 신규 설정된 저당대부를 인수하는 조건으로 매매한 것이다. \par
- 거래금액은 10억원이며, 대부비율은 70\%여서 차액인 30\%만이 매도인에게 지불되었다. \par
- 저당기간은 20년이고, 저당이자율은 연간 10\%이며, 현재의 전형적인 시장이자율은 연간 12\%이다. \par
- 한편, 저당대부는 만기까지 존속되는 것을 가정하고, 원리금의 지불은 편의상 연단위로 계산하며, < 자료 4 >의 연복리표를 활용한다. \par
< 자료 4 > 연복리표 (20년, 소수점 다섯째 자리에서 절사함) \par
{\small\centering\renewcommand{\arraystretch}{1.25}
\resizebox{\ifdim\width>\linewidth\linewidth\else\width\fi}{!}{
\begin{tabular}{|c|c|c|c|c|c|c|}
\hline
\makecell{이자율} & \makecell{일시불\\내가계수} & \makecell{연금\\내가계수} & \makecell{감채기금\\계수} & \makecell{일시불\\현가계수} & \makecell{연금\\현가계수} & \makecell{저당상수} \\
\hline
\makecell{10\%} & \makecell{6.7275} & \makecell{57.2749} & \makecell{0.0174} & \makecell{0.1486} & \makecell{8.5135} & \makecell{0.1174} \\
\hline
\makecell{12\%} & \makecell{9.6462} & \makecell{72.0524} & \makecell{0.0138} & \makecell{0.1036} & \makecell{7.4694} & \makecell{0.1338} \\
\hline
\end{tabular}}\par}

< 자료 5 > 낙찰사례(3)에 대한 개요 \par
1. 토지: 면적은 400m\textsuperscript{2}이고, 별도 명도비용의 발생이 예상된다. \par
2. 건물: 없음 \par
3. 경매 관련사항 \par
- 낙찰사례(3)의 가격은 최초의 경매개시 가격의 50\%에서 결정되었으며, 평가선례(4)는 낙찰사례 자체를 대상으로 한 경매목적의 평가로 평가액은 10억원이었고, 점유자 및 적치물은 고려치 못한 상태에서 평가가 이루진 것으로 조사된다. \par
- 낙찰자는 점유자 및 그 적치물의 상태를 고려(인수조건)하여 입찰에 참여하였고, 이에 따른 명도비용은 약 1억5천만원이 소요될 것으로 예상되며, 기간이자 및 자가노력비는 고려치 아니하였다. \par
- 관련 통계에 의하면, 사례가 속한 지역의 해당기간 평균적인 경매의 낙찰률(최초 평가가격 대비 낙찰가격)은 다음과 같이 조사된다. \par
{\small\centering\renewcommand{\arraystretch}{1.25}
\resizebox{\ifdim\width>\linewidth\linewidth\else\width\fi}{!}{
\begin{tabular}{|c|c|c|c|c|}
\hline
\makecell{구분} & \makecell{상업지역} & \makecell{주거지역} & \makecell{녹지지역} & \makecell{지역평균} \\
\hline
\makecell{낙찰률} & \makecell{90\%} & \makecell{80\%} & \makecell{70\%} & \makecell{80\%} \\
\hline
\end{tabular}}\par}

- 한편, 낙찰률의 적용은 낙찰사례가 속한 지역이 주거지역이지만, 녹지지역과 연접하기 때문에 녹지지역과 주거지역 낙찰률의 평균을 적용하였다. \par
\end{longquestion}
